\documentclass[aspectratio=169]{beamer}
\usetheme{Madrid}
\usepackage{amsmath, amssymb}

\title{Matrix Representations of Linear Transformations}
\subtitle{Calm, Clear, and Actually Understandable}
\author{MATH 304 -- Rewritten for Sanity}
\date{}

\begin{document}

%------------------------------------------------
\begin{frame}
\titlepage
\end{frame}

%------------------------------------------------
\begin{frame}{First: Breathe. This Is Not Magic.}
You are not bad at math.  

You were shown symbols without meaning.

Today we fix that.

\bigskip

Core idea of this lecture:

\[
\textbf{Every linear transformation is just a matrix in disguise.}
\]

\bigskip

If you understand:
\begin{itemize}
\item what a basis is,
\item what coordinates mean,
\item and what columns represent,
\end{itemize}

then this entire topic collapses into a pattern.
\end{frame}

%------------------------------------------------
\begin{frame}{CHEAT SHEET (1/3)}
\begin{columns}[t]
\column{0.48\textwidth}
\textbf{Definition: Linear Transformation}

A map \( L: V \to W \) is linear if

\[
L(\alpha v_1 + \beta v_2)
=
\alpha L(v_1) + \beta L(v_2)
\]

for:
\begin{itemize}
\item vectors \(v_1, v_2 \in V\)
\item scalars \(\alpha, \beta \in \mathbb{R}\)
\end{itemize}

\bigskip

\textbf{Standard Basis in \( \mathbb{R}^n \)}

\[
e_1 =
\begin{pmatrix}
1\\0\\\vdots\\0
\end{pmatrix},
\quad
e_2 =
\begin{pmatrix}
0\\1\\\vdots\\0
\end{pmatrix}
\]

\column{0.48\textwidth}
\textbf{Big Theorem (Standard Case)}

If \(L: \mathbb{R}^n \to \mathbb{R}^m\) is linear,

then there exists an \(m \times n\) matrix \(A\) such that:

\[
L(x) = Ax
\]

Columns of \(A\) are:

\[
a_j = L(e_j)
\]

\bigskip

\textbf{Translation:}

To find the matrix:

\begin{enumerate}
\item Plug in each basis vector.
\item Those outputs become columns.
\end{enumerate}

That is it.
\end{columns}
\end{frame}

%------------------------------------------------
\begin{frame}{CHEAT SHEET (2/3)}
\begin{columns}[t]
\column{0.48\textwidth}
\textbf{Coordinates Relative to a Basis}

If \(E = \{v_1, \dots, v_n\}\),

and

\[
v = x_1 v_1 + \dots + x_n v_n
\]

then

\[
[v]_E =
\begin{pmatrix}
x_1\\
\vdots\\
x_n
\end{pmatrix}
\]

These numbers are just weights.

\column{0.48\textwidth}
\textbf{Matrix Relative to Bases}

If
\[
L: V \to W
\]

with bases
\[
E = \{v_1,\dots,v_n\},\quad
F = \{w_1,\dots,w_m\}
\]

then

\[
[L(v)]_F = A [v]_E
\]

Columns of \(A\):

\[
a_j = [L(v_j)]_F
\]

\textbf{Pattern:}

Input basis $\rightarrow$ apply $L$ $\rightarrow$ rewrite in output basis.
\end{columns}
\end{frame}

%------------------------------------------------
\begin{frame}{CHEAT SHEET (3/3)}
\begin{columns}[t]
\column{0.48\textwidth}
\textbf{Change of Basis Trick}

If \(F\) is not standard,

let

\[
B = (b_1,\dots,b_m)
\]

Then:

\[
a_j = B^{-1} L(v_j)
\]

\bigskip

\textbf{Row Reduction Shortcut}

The RREF of

\[
(b_1,\dots,b_m \mid L(v_1),\dots,L(v_n))
\]

becomes

\[
(I \mid A)
\]

\column{0.48\textwidth}
\textbf{Polynomial Case}

Derivative map:

\[
D(p) = p'
\]

To build matrix:

\begin{enumerate}
\item Differentiate each basis polynomial.
\item Rewrite result in output basis.
\item Those coordinate vectors form columns.
\end{enumerate}

\bigskip

Everything follows this recipe.
\end{columns}
\end{frame}

%------------------------------------------------
\begin{frame}{Why This Works (Intuition)}
A linear transformation respects addition and scaling.

So if

\[
x = x_1 e_1 + \dots + x_n e_n
\]

then linearity forces:

\[
L(x) = x_1 L(e_1) + \dots + x_n L(e_n)
\]

\bigskip

That means:

\[
L(x) = (L(e_1) \ \dots \ L(e_n))
\begin{pmatrix}
x_1\\
\vdots\\
x_n
\end{pmatrix}
\]

\bigskip

That matrix is \(A\).

This is not a coincidence.

It is unavoidable.
\end{frame}

%------------------------------------------------
\begin{frame}{Example 1 (Part 1)}
Let

\[
L: \mathbb{R}^3 \to \mathbb{R}^2
\]

\[
L(x) = (x_1 + x_2,\ x_2 + x_3)
\]

Define:

\[
x =
\begin{pmatrix}
x_1\\
x_2\\
x_3
\end{pmatrix}
\]

Goal: Find matrix \(A\) so \(L(x)=Ax\).

\end{frame}

%------------------------------------------------
\begin{frame}{Example 1 (Part 2)}
Standard basis in \( \mathbb{R}^3 \):

\[
e_1 =
\begin{pmatrix}1\\0\\0\end{pmatrix},
e_2 =
\begin{pmatrix}0\\1\\0\end{pmatrix},
e_3 =
\begin{pmatrix}0\\0\\1\end{pmatrix}
\]

Compute:

\[
L(e_1) = (1,0)
\]

\[
L(e_2) = (1,1)
\]

\[
L(e_3) = (0,1)
\]

\end{frame}

%------------------------------------------------
\begin{frame}{Example 1 (Final)}
Columns are outputs:

\[
A =
\begin{pmatrix}
1 & 1 & 0\\
0 & 1 & 1
\end{pmatrix}
\]

WHY?

Because:

\[
Ax =
x_1
\begin{pmatrix}1\\0\end{pmatrix}
+
x_2
\begin{pmatrix}1\\1\end{pmatrix}
+
x_3
\begin{pmatrix}0\\1\end{pmatrix}
\]

which equals

\[
(x_1 + x_2,\ x_2 + x_3)
\]

It matches perfectly.
\end{frame}

%------------------------------------------------
\begin{frame}{Polynomial Example (Part 1)}
Let

\[
D: P_2 \to P_1
\]

\[
D(p) = p'
\]

Basis:

\[
E = [x^2, x, 1]
\]

\[
F = [x, 1]
\]

\end{frame}

%------------------------------------------------
\begin{frame}{Polynomial Example (Part 2)}
Differentiate basis elements:

\[
D(x^2) = 2x
\]

Rewrite:

\[
2x = 2\cdot x + 0\cdot 1
\]

So

\[
[D(x^2)]_F =
\begin{pmatrix}
2\\0
\end{pmatrix}
\]

\end{frame}

%------------------------------------------------
\begin{frame}{Polynomial Example (Part 3)}
\[
D(x) = 1 = 0\cdot x + 1\cdot 1
\]

\[
[D(x)]_F =
\begin{pmatrix}
0\\1
\end{pmatrix}
\]

\[
D(1)=0
\]

\[
[D(1)]_F =
\begin{pmatrix}
0\\0
\end{pmatrix}
\]

\end{frame}

%------------------------------------------------
\begin{frame}{Polynomial Example (Final)}
Matrix:

\[
A =
\begin{pmatrix}
2 & 0 & 0\\
0 & 1 & 0
\end{pmatrix}
\]

If

\[
p(x)=ax^2+bx+c
\]

then

\[
[p]_E =
\begin{pmatrix}
a\\b\\c
\end{pmatrix}
\]

\[
[D(p)]_F =
A
\begin{pmatrix}
a\\b\\c
\end{pmatrix}
=
\begin{pmatrix}
2a\\b
\end{pmatrix}
\]

Which means

\[
D(p)=2ax+b
\]

It works exactly as expected.
\end{frame}

%------------------------------------------------
\begin{frame}{Homework Survival Mindset}
Every problem reduces to:

\begin{enumerate}
\item Identify input basis.
\item Apply transformation to each basis vector.
\item Rewrite outputs in output basis.
\item Those coordinate vectors are columns.
\end{enumerate}

If change of basis appears:

\[
a_j = B^{-1}L(v_j)
\]

\bigskip

There is no hidden trick.

Only repeated pattern.
\end{frame}

%------------------------------------------------
\begin{frame}{Extra Example (Deep WHY)}
Let

\[
L: \mathbb{R}^2 \to \mathbb{R}^2
\]

\[
L(x_1,x_2) = (2x_1 - x_2,\ 3x_2)
\]

Step 1:

\[
L(e_1)=(2,0)
\]

\[
L(e_2)=(-1,3)
\]

\end{frame}

%------------------------------------------------
\begin{frame}{Extra Example (Solution)}
Matrix:

\[
A=
\begin{pmatrix}
2 & -1\\
0 & 3
\end{pmatrix}
\]

WHY this works:

\[
Ax=
x_1
\begin{pmatrix}2\\0\end{pmatrix}
+
x_2
\begin{pmatrix}-1\\3\end{pmatrix}
\]

\[
=
(2x_1-x_2,\ 3x_2)
\]

Same structure again.

Different numbers.

Same pattern.
\end{frame}

%------------------------------------------------
\begin{frame}{Final Words}
This topic is not advanced.

It is repetitive.

\bigskip

Linear transformation
=
Matrix
+
Choice of basis.

\bigskip

You were overwhelmed by notation.

But the structure is simple.

And you can absolutely do this.
\end{frame}

\end{document}